\scanpage{96}

en faisant attention que pour $i$ et $j$ incomparables et $X$ métrisable, on \deleted{\unsure{trouve}} aura $\mathcal V_{ij}=\varnothing$ \deleted{\unsure{et}} (donc pas la peine d'en parler!), et pour $i=j$, on aura $\mathcal V_{ij}=\mathcal V_{ii}=\mathcal V_i$ (donc pas la peine de l'introduire, $\mathcal V_i$ suffit!), enfin \deleted{\unsure{et}} si $i$ et $j$ sont distincts et comparables, \deleted{\unsure{donc}} comme $\mathcal V_{ij}=\mathcal V_{ji}$, il suffit d'introduire $\mathcal V_{ij}$ pour \deleted{\unsure{les}} $i<j$. Plus généralement, on introduit les
\[(98)\qquad\mathcal V_{i_0\ldots i_n}=\mathcal V_{i_0}\cap\cdots\cap\mathcal V_{i_n},\]
mais on se borne en règle générale au cas où $\{i_0,\ldots,i_n\}$ est un drapeau (sinon on trouve $\varnothing$) et où en fait ils sont rangés en ordre str. croissant
\[i_0<i_1<\cdots<i_n\]
(pour éviter des redondances). Les protopos $^*$ sont les ``pierres de construction'' élémentaires, pour en construire des tas d'autres, par \uline{$\lim$} (finies si $I$ fini), compte tenu des morphismes semi-simpliciaux évidents entre les $\mathcal V_d$ $(d=(i_0,\ldots,i_n)\in\mathrm{Drap}(I))$.
\margin{$^*$(et les \unsure{protopos} \unsure{homotopes} $\mathcal V_{i_0\ldots i_n}^{\unsure{Z}}$, \unsure{cf.} plus bas)}
\margin{95}
